\documentclass[11pt]{article}

\usepackage[margin=1in]{geometry}
\usepackage{amsmath,amssymb,amsthm,mathtools}
\usepackage{enumitem}
\usepackage{booktabs}
\usepackage{tabularx}
\usepackage{microtype}
\usepackage[hidelinks]{hyperref}
\usepackage{xurl}

\newtheorem{definition}{Definition}
\newtheorem{lemma}{Lemma}
\newtheorem{remark}{Remark}

\title{Stable Release Spine Contracts for Successor Maintenance:\\[0.5ex]
\large Stage Ownership, Handoff Rules, and Auditable Walk-Throughs}
\author{}
\date{Unified archive note v0.14 (2026-03-19; archive v400)}

\begin{document}
\maketitle

\begin{abstract}
Later papers need a stable way to say how one released claim evolved without re-telling the compare, replay, obligation, closure, and downstream-sentence story each time.
This note gives that layer one home.
It exports a checked release spine
\emph{diff $\to$ classify $\to$ obligate $\to$ close $\to$ summarize $\to$ support $\to$ cite $\to$ quote $\to$ clause},
together with one tiny contract for each stage: what inputs it may read, what output family it owns, what kind of assertion it is allowed to support, and when later notes must cite that stage directly rather than let wrapper prose drift across boundaries.
The result is a stable publication backbone for release evolution that remains auditable and non-decorative.
It also gives later papers a principled answer to a practical question: what is the earliest release stage that is actually sufficient for the claim being cited?
\end{abstract}

\paragraph{Non-redundancy note.}
Synthesis~18 owns drift classes, Synthesis~20 owns replay mechanics, Synthesis~33 owns the role-by-role obligation/closure bridge inside the obligate$\to$close handoff, Synthesis~34 owns the narrow summary-stage status-envelope / lineage-notice split, and Synthesis~31 owns the downstream suffix once the release spine reaches support/citation/quote/clause reuse. Synthesis~45--54 then own the answer-side companion objects that may sit beside that suffix without becoming new release stages, and Synthesis~74--75 own the paired terminal reuse / cutover discipline that later terse papers apply once those checked-answer owners already exist.\cite{series:challengeanswercapsules,series:challengeanswersheets,series:challengeanswercards,series:challengeanswercarrierslots,series:challengeansweraudittrails,series:challengeanswercitationdockets,series:challengeanswerpublicationprofiles,series:challengeanswerdisclosurepackets,series:challengeanswerescalationladders,series:challengeanswerreviewmenus,series:pairedterminalcitationdiscipline,series:pairedterminalcutoverrules}
This note owns only the stable stage vocabulary and its handoff discipline.

\paragraph{Scope.}
This is not a new mechanism note and not a new receipt primitive.
It is a publication-interface note for successor maintenance after a claim has already been released.

\paragraph{Imports.}
Contract-object vocabulary is imported from Synthesis~0.\cite{series:contractobjects}
Plan/state drift language is imported from Synthesis~18.\cite{series:planstateequiv}
Replay verdict structure is imported from Synthesis~20.\cite{series:auditorreplay}
The downstream normalization suffix is imported from Synthesis~31.\cite{series:downstreamnormalforms}
The obligate/close bridge is imported from Synthesis~33.\cite{series:releaseclosure}
The summary-stage status-envelope split is imported from Synthesis~34.\cite{series:statusenvelopes}
The answer-side terminal reuse / cutover discipline is imported from Synthesis~54 and Synthesis~74--75, the concrete worked instantiation is imported from Synthesis~17, the cross-series theorem-to-review bridge is imported from Synthesis~55, and the source-only public-request carryforward bridge that reuses this release spine across successor cuts is imported from Synthesis~60.\cite{series:challengeanswerreviewmenus,series:pairedterminalcitationdiscipline,series:pairedterminalcutoverrules,series:workedexample,series:seriesspines,series:publicrequestcarryforward}

\paragraph{Exports.}
This note exports (i) a release-stage handoff contract, (ii) a validity rule for a stable release spine, (iii) a stage-locality rule that separates owned outputs from later wrappers, (iv) a shared-stage-vocabulary rule saying that later routes, series spines, and review-side mirrors may repeat stage labels only as imported release-spine vocabulary, (v) a stage-walkthrough witness that turns the spine into one auditable example rather than a decorative pipeline diagram, and (vi) a local-reopen rule saying that later terse papers should reopen this note only when the exact release-stage family, minimal stage cut, or stage-local handoff boundary is itself live. Otherwise, once a question has already been reduced to a checked-answer handoff, later terse papers should default to the answer-side terminal citation normal form in Synthesis~54 together with the paired terminal discipline and paired cutover rule in Synthesis~74--75. This note does not try to own the internal role rows used by the obligate$\to$close bridge, the answer-side companion objects that sit beside the downstream suffix, or the later source-only public-request carryforward tokens that say whether a predecessor-complete promise stays complete or reopens under the successor release.\cite{series:challengeanswerreviewmenus,series:pairedterminalcitationdiscipline,series:pairedterminalcutoverrules}

\section{Stable stage ownership}

\begin{definition}[Release-stage handoff contract]
For one stage token $z$ in the release spine, a \emph{release-stage handoff contract}
\[
H(z)=(I_z,O_z,F_z,A_z,R_z)
\]
consists of a finite input-artifact list $I_z$, one canonical output artifact $O_z$, one declared output field family $F_z$ owned by that stage, one assertion scope $A_z$ describing what later notes may legitimately claim from that stage, and one handoff rule $R_z$ saying when a later note must cite that stage directly versus continue leftward to a narrower object.
The contract does not replace the stage output.
Its job is only to stop stage order from becoming a decorative flowchart with ambiguous ownership of numeric deltas, decision tokens, or reusable sentence text.
\end{definition}

\begin{definition}[Stable release spine]
A \emph{stable release spine} is an ordered list of stage contracts
\[
\Sigma=(H_1,\dots,H_9)
\]
whose stage labels are exactly
\[
\textsf{diff},\textsf{classify},\textsf{obligate},\textsf{close},\textsf{summarize},\textsf{support},\textsf{cite},\textsf{quote},\textsf{clause},
\]
and whose inputs are monotone: each stage may read earlier released maintenance objects, but no later stage may silently re-own a field family already declared by an earlier stage.
\end{definition}

\begin{remark}[Stage-locality rule]
Later stages may compress earlier results for readability, but they do not replace them.
A summary sentence may quote a continuity token; it does not own that token.
A clause pack may compress a numeric delta; it does not own the delta.
This is the single rule that keeps the spine auditable.
\end{remark}

\begin{remark}[Shared-stage-vocabulary rule]
Later routes, review menus, series spines, and worked examples may mirror repeated stage labels such as \textsf{diff} or \textsf{close}, but only as imported release-spine vocabulary.
The release spine still owns the stage order, the per-stage stop rules, and the handoff contract attached to each such label.
So a later bridge may say that an answer depends on the \textsf{classify} and \textsf{support} stages without becoming the owner of either stage policy.
\end{remark}

\begin{table}[t]
\centering
\footnotesize
\begin{tabularx}{\linewidth}{@{}l >{\raggedright\arraybackslash}p{0.20\linewidth} >{\raggedright\arraybackslash}p{0.26\linewidth} >{\raggedright\arraybackslash}X@{}}
\toprule
Stage & Canonical output & Owned field family & Earliest stop question \\
\midrule
Diff & compare report & observed diffs; winning class & What changed? \\
Classify & continuity verdict & continuity decision; changed slice; replay scope & What replay scope and continuity class follow? \\
Obligate & obligation profile (+ action matrix) & required outputs by action class & Which public outputs must refresh? \\
Close & closure verdict & closure status; bound shipped outputs & Were the required outputs shipped? \\
Summarize & lineage notice (+ status envelope) & outward-facing status sentence; narrow public-status token & What public successor-status answer is now valid? \\
Support & claim-support map & clause-to-evidence bindings & Which narrower object supports the public clause? \\
Cite & citation map & smallest-sufficient citation choices & What is the smallest sufficient citation target? \\
Quote & quote map & exact field selections & Which exact field should later notes quote? \\
Clause & clause pack & reusable short wording & Which checked short clause is reusable? \\
\bottomrule
\end{tabularx}
\caption{The stable release spine owns stage vocabulary and field-family boundaries, not new mechanism theory.}
\label{tab:spine}
\end{table}

\section{Minimal downstream stop rules}

\begin{definition}[Minimal stage cut]
Fix one downstream question $q$ about a maintained base$\to$successor package. The \emph{minimal stage cut} $c(q)$ is the earliest stage in the stable release spine whose owned field family is sufficient to answer $q$ without appealing to later wrapper wording.
\end{definition}

\begin{remark}[Earliest sufficient stop]
The minimal stage cut is the operational payoff of the spine. If a later note asks only what changed, it should stop at \textsf{diff}. If it asks what had to be replayed, it should stop at \textsf{classify}. If it asks whether the package is publication-complete, it must reach \textsf{close}. If it asks what short sentence readers saw, it may stop at \textsf{summarize}; and if it asks which exact reusable clause is already checked, it may continue to \textsf{clause}.
\end{remark}

\paragraph{Boundary of this rule.} This rule is deliberately release-local. If the question still needs the theorem root or the receipt-facing claim surface, the series spine in Synthesis~55 composes this stage cut with those outer owners instead of widening the release spine itself. If the question has already reached a checked-answer handoff, later terse papers should then default to the answer-side terminal citation normal form in Synthesis~54 together with the paired terminal discipline and paired cutover rule in Synthesis~74--75 rather than re-citing stage labels from this note.\cite{series:seriesspines,series:challengeanswerreviewmenus,series:pairedterminalcitationdiscipline,series:pairedterminalcutoverrules}

\paragraph{Later-terse-paper citation posture.} Later terse papers should cite Synthesis~32 only when the live issue is the release-stage order itself, the earliest sufficient stage cut for a release-evolution question, or the exact handoff boundary between adjacent release stages. They should stop at Synthesis~33 when the live issue is the obligation-profile row set or closure-ledger satisfaction verdict, stop at Synthesis~34 when the live issue is the narrow status token or its outward-facing lineage wrapper, and stop at Synthesis~31 when the release question has already reduced to downstream summary / support / citation / quote / clause reuse. Once a checked-answer handoff already exists, later terse papers should widen no further through this note and instead default to the paired terminal citation discipline and paired cutover rule in Synthesis~74--75.\cite{series:releaseclosure,series:statusenvelopes,series:downstreamnormalforms,series:pairedterminalcitationdiscipline,series:pairedterminalcutoverrules}

\begin{table}[t]
\centering
\small
\begin{tabularx}{\linewidth}{@{}l l X@{}}
\toprule
Downstream question & Minimal stage cut & Why later wrapper stages are unnecessary \\
\midrule
What changed between base and successor? & \textsf{diff} & The compare report already owns the observed-diff family and winning diff class. \\
Which slice had to be replayed or re-certified? & \textsf{classify} & The continuity verdict already owns replay scope, changed slice, and continuity class. \\
Which public outputs had to be refreshed, and were they shipped? & \textsf{obligate}/\textsf{close} & The obligation profile and closure verdict already answer duty and completion without later prose. \\
What exact outward-facing successor sentence was shown? & \textsf{summarize} & The lineage notice and status envelope already own that emitted summary-layer answer. \\
Which short reusable wording is already checked for later papers? & \textsf{clause} & Only the clause pack owns the final reusable wording; earlier stages own narrower facts, not the checked clause. \\
\bottomrule
\end{tabularx}
\caption{The stable spine is most useful when downstream papers stop at the earliest sufficient stage rather than re-telling the whole maintenance chain.}
\label{tab:minimal-stage-cut}
\end{table}

\section{What the contracts prevent}
Without stage contracts, three failure modes recur:
\begin{enumerate}[leftmargin=*]
\item a quote-stage token is cited as if it were the continuity decision itself;
\item a summary sentence is cited as if it owned numeric deltas that actually belong to the compare report; and
\item a clause-pack sentence is reused as if it were portable across release pairs merely because the wording template still reads plausibly.
\end{enumerate}
The handoff contract blocks all three by naming one owner for each field family and one explicit rule for when later notes must stop compressing and cite the owner directly.

\begin{lemma}[Stage-local audit boundary]
Assume a stable release spine $\Sigma$ and suppose every later-stage assertion obeys the handoff rule of the stage whose owned field family it uses.
Then no asserted token, numeric delta, or reusable sentence needs an implicit ownership inference: every challenged claim resolves either to the owning stage contract or to a narrower object explicitly named by that contract.
\end{lemma}

\begin{proof}
Take any challenged downstream statement.
If it is only a wrapper sentence, the relevant handoff rule may stop at the summarize, support, cite, quote, or clause stage.
If it asserts a token or numeric field, the stage-locality rule forbids that statement from terminating at a later wrapper stage and sends the audit path back to the owner named in the relevant contract.
Because every stage contract fixes its output family and because later stages may compress but not re-own earlier families, ownership never depends on prose interpretation alone.
\end{proof}

\section{Auditable walk-throughs}

\begin{definition}[Stage-walkthrough witness]
For one concrete base$\to$successor comparison, a \emph{stage-walkthrough witness}
\[
W=(w_1,\dots,w_9)
\]
is an ordered list with one entry per spine stage.
Each entry names the stage label, the concrete output artifact, the owned field family realized in that artifact, the concrete inputs read for that step, and one short explanation of why later notes must not let that field family drift into a later wrapper.
\end{definition}

\begin{remark}[Why a walk-through matters]
A stage list alone is too easy to wave at.
A stage-walkthrough witness turns the spine into one auditable worked example: it shows that the stage order is actually instantiated by concrete objects and that the ownership boundary is checkable field-family by field-family.
\end{remark}

\section{Worked example pointer}
The maintained worked bundle ships two complementary objects.\cite{series:workedexample}
The first, \texttt{example\_release\_spine.json}, freezes the stable stage order together with one contract per stage.
The second, \texttt{example\_release\_stage\_walkthrough.json}, instantiates those contracts for the canned Case-B successor by listing the concrete output artifact, owned field family, resolved inputs, and downstream citation boundary for each stage.
Together they give later papers one non-decorative release-evolution backbone to cite, while Synthesis~31 remains the home for sentence-level reuse once the spine reaches its downstream suffix. The answer-side companion stack may then be assembled beside that suffix for particular audiences and questions, but it does not change the stage order itself; once that checked-answer handoff exists, later terse papers should default to the answer-side terminal citation normal form and paired cutover rule rather than restating release-stage vocabulary locally.\cite{series:challengeanswercapsules,series:challengeanswersheets,series:challengeanswercards,series:challengeanswercarrierslots,series:challengeansweraudittrails,series:challengeanswercitationdockets,series:challengeanswerpublicationprofiles,series:challengeanswerdisclosurepackets,series:challengeanswerescalationladders,series:challengeanswerreviewmenus,series:pairedterminalcitationdiscipline,series:pairedterminalcutoverrules}

\paragraph{Why this helps the series.}
The archive wants short papers without hidden maintenance drift.
Owning the release spine here lets the worked example stay crisp, lets downstream-normal-form notes stay focused on sentence reuse, and gives later papers one stable answer to ``how did this release evolve?'' without re-narrating the whole maintenance stack.

\begin{thebibliography}{99}\small

\bibitem{series:contractobjects}
Anonymous.
\newblock Contract Objects for Anonymous-DHT Claims: Commitments, Menus, Receipts, and Release Bindings.
\newblock Series synthesis note (Synthesis~0), 2026. (This archive.)

\bibitem{series:planstateequiv}
Anonymous.
\newblock Plan and State Equivalence: Digest-Bound Replay Hooks and Drift Taxonomy.
\newblock Series synthesis note (Synthesis~18), 2026. (This archive.)

\bibitem{series:auditorreplay}
Anonymous.
\newblock Auditor Replay Recipe for Anonymous-DHT Receipts: A Minimal Checklist and Verdict Tuple.
\newblock Series synthesis note (Synthesis~20), 2026. (This archive.)

\bibitem{series:downstreamnormalforms}
Anonymous.
\newblock Downstream Normal Forms for Successor Maintenance: Summary, Support, Citation, Quote, and Clause Discipline.
\newblock Series synthesis note (Synthesis~31), 2026. (This archive.)

\bibitem{series:releaseclosure}
Anonymous.
\newblock Release Obligations and Closure Ledgers for Successor Maintenance: Required Roles, Row Satisfaction, and Completion Tokens.
\newblock Series synthesis note (Synthesis~33), 2026. (This archive.)

\bibitem{series:statusenvelopes}
Anonymous.
\newblock Successor Status Envelopes and Lineage Notices for Release Maintenance: Narrow Status Tokens, Reader Pointer Order, and Public Answers.
\newblock Series synthesis note (Synthesis~34), 2026. (This archive.)

\bibitem{series:challengeanswercapsules}
Anonymous.
\newblock Challenge Answer Capsules for Successor Maintenance: Audience-Question Minimal Answer Bundles from Certified Piece Hooks.
\newblock Series synthesis note (Synthesis~45), 2026. (This archive.)

\bibitem{series:challengeanswersheets}
Anonymous.
\newblock Challenge Answer Sheets for Successor Maintenance: Audience-Question Piece-Complete Handoffs from Answer Capsules.
\newblock Series synthesis note (Synthesis~46), 2026. (This archive.)

\bibitem{series:challengeanswercards}
Anonymous.
\newblock Challenge Answer Cards for Successor Maintenance: Audience-Question Released Answers from Answer Sheets and Sentence Locks.
\newblock Series synthesis note (Synthesis~47), 2026. (This archive.)

\bibitem{series:challengeanswercarrierslots}
Anonymous.
\newblock Challenge Answer Carrier Slots for Successor Maintenance: Exact Shipped Fields and Reusable Clauses from Answer Cards.
\newblock Series synthesis note (Synthesis~48), 2026. (This archive.)

\bibitem{series:challengeansweraudittrails}
Anonymous.
\newblock Challenge Answer Audit Trails for Successor Maintenance: Leftward Reader Walks from Carrier Slots to Narrow Terminals.
\newblock Series synthesis note (Synthesis~49), 2026. (This archive.)

\bibitem{series:challengeanswercitationdockets}
Anonymous.
\newblock Challenge Answer Citation Dockets for Successor Maintenance: Brief Provenance Budgets from Audit Trails.
\newblock Series synthesis note (Synthesis~50), 2026. (This archive.)

\bibitem{series:challengeanswerpublicationprofiles}
Anonymous.
\newblock Challenge Answer Publication Profiles for Successor Maintenance: Default Public vs Requestable Review Surfaces from Citation Dockets.
\newblock Series synthesis note (Synthesis~51), 2026. (This archive.)

\bibitem{series:challengeanswerdisclosurepackets}
Anonymous.
\newblock Challenge Answer Disclosure Packets for Successor Maintenance: Exact Requestable Review Bundles from Publication Profiles.
\newblock Series synthesis note (Synthesis~52), 2026. (This archive.)

\bibitem{series:challengeanswerescalationladders}
Anonymous.
\newblock Challenge Answer Escalation Ladders for Successor Maintenance: Ordered Brief-to-Audit Reveal Policies from Publication Profiles and Disclosure Packets.
\newblock Series synthesis note (Synthesis~53), 2026. (This archive.)

\bibitem{series:challengeanswerreviewmenus}
Anonymous.
\newblock Challenge Answer Review Menus for Successor Maintenance: Smallest-Sufficient Referee Handoffs from Dockets, Profiles, Packets, and Ladders.
\newblock Series synthesis note (Synthesis~54), 2026. (This archive.)

\bibitem{series:pairedterminalcitationdiscipline}
Anonymous.
\newblock Paired Terminal Citation Discipline for Review Spines: Stable Answer Stops and Adjacent Request Widening.
\newblock Series synthesis note (Synthesis~74), 2026. (This archive.)

\bibitem{series:pairedterminalcutoverrules}
Anonymous.
\newblock Paired Terminal Cutover Rules for Review Spines: When Later Terse Papers Stop, Widen, or Reopen.
\newblock Series synthesis note (Synthesis~75), 2026. (This archive.)

\bibitem{series:seriesspines}
Anonymous.
\newblock Series Spines for Anonymous-DHT Receipts: Theorem Roots, Checked Routes, Request Bridges, and Stable Reuse Across Release Evolution.
\newblock Series synthesis note (Synthesis~55), 2026. (This archive.)

\bibitem{series:publicrequestcarryforward}
Anonymous.
\newblock Public Request Carryforward for Source-Only Successor Releases: When Paper-Level Request Status Persists, Advances, or Reopens across Release Evolution.
\newblock Series synthesis note (Synthesis~60), 2026. (This archive.)

\bibitem{series:workedexample}
Anonymous.
\newblock Worked Example: Receipt Interlock for an Anonymous-DHT Reader-Privacy Claim.
\newblock Series synthesis note (Synthesis~17), 2026. (This archive.)

\end{thebibliography}
\end{document}
